\chapter{CƠ SỞ LÝ THUYẾT VÀ PHƯƠNG PHÁP THỰC HIỆN}

\section{Tổng quan dữ liệu World Happiness Report}

\subsection{Bản chất của dữ liệu}

Dữ liệu của đề tài được lưu dưới dạng tệp CSV theo từng năm. Mỗi dòng dữ liệu mô tả thông tin của một quốc gia trong một năm báo cáo. Cấu trúc dữ liệu gồm tên quốc gia, khu vực, thứ hạng hạnh phúc, điểm hạnh phúc và các yếu tố liên quan.

Điểm hạnh phúc được sử dụng làm biến mục tiêu để xếp hạng và so sánh mức độ hạnh phúc giữa các quốc gia. Thứ hạng hạnh phúc cho biết vị trí tương đối của quốc gia trong từng năm. GDP bình quân đầu người, hỗ trợ xã hội, tuổi thọ khỏe mạnh, tự do lựa chọn, sự hào phóng và nhận thức về tham nhũng là các biến giải thích được khảo sát trong hệ thống.

\subsection{Cấu trúc dữ liệu giai đoạn 2015--2019}

\begin{table}[H]
\centering
\begin{tabular}{|p{4cm}|p{8.5cm}|}
\hline
\textbf{Nhóm dữ liệu} & \textbf{Tên trường/chức năng} \\
\hline
Thông tin định danh & Country, Region, Year \\
\hline
Chỉ số hạnh phúc & Happiness Rank, Happiness Score \\
\hline
Yếu tố kinh tế & GDP per capita \\
\hline
Yếu tố xã hội & Social support, Life Expectancy, Freedom \\
\hline
Yếu tố hành vi và nhận thức & Generosity, Perceptions of corruption \\
\hline
\end{tabular}
\caption{Cấu trúc dữ liệu chính của dự án}
\label{tab:co-cau-du-lieu}
\end{table}

Bảng \ref{tab:co-cau-du-lieu} liệt kê các trường được sử dụng thống nhất trong dashboard. Một số tệp gốc còn chứa các cột bổ sung như khoảng tin cậy, độ dư Dystopia hoặc sai số chuẩn. Vì các cột này không có đầy đủ ở mọi năm nên không được sử dụng làm chỉ số phân tích chung.

\section{Tiền xử lý dữ liệu}

\subsection{Chuẩn hóa cấu trúc}

Do dữ liệu của các năm được xuất từ nhiều nguồn tương tự nhưng có tên cột không đồng nhất, hệ thống ánh xạ các tên trường về tên chung. Một số ánh xạ tiêu biểu là:

\begin{itemize}
    \item \textit{Country or region} được chuyển thành \textit{Country}.
    \item \textit{Overall rank} hoặc \textit{Happiness.Rank} được chuyển thành \textit{Happiness Rank}.
    \item \textit{Score} hoặc \textit{Happiness.Score} được chuyển thành \textit{Happiness Score}.
    \item Các tên GDP, tuổi thọ, hỗ trợ xã hội và nhận thức tham nhũng của từng năm được chuẩn hóa theo tên chung.
\end{itemize}

\subsection{Bổ sung khu vực}

Mỗi bản ghi được gắn với khu vực địa lý để hỗ trợ so sánh vùng. Hệ thống sử dụng dữ liệu khu vực của hai năm tham chiếu 2015 và 2016 để xây dựng ánh xạ quốc gia -- khu vực. Với các tên quốc gia khác nhau giữa các năm, module tiền xử lý định nghĩa một số bí danh như Hong Kong, Taiwan, North Macedonia và Trinidad and Tobago.

\subsection{Xử lý giá trị thiếu}

Các cột số được chuyển sang kiểu số; giá trị không hợp lệ sẽ được xử lý bằng \textit{median} của cột trong từng tệp dữ liệu. Cách làm này giúp hệ thống giữ kích thước mẫu và không loại toàn bộ bản ghi chỉ vì một số ít giá trị thiếu.

Khi tải dữ liệu cho dashboard, \textit{dashboard\_utils.py} tiếp tục chuyển các trường cần phân tích sang kiểu số bằng \textit{pd.to\_numeric(..., errors="coerce")}. Sau đó, dữ liệu trùng quốc gia trong cùng năm được loại bỏ và kết quả được sắp xếp theo năm và thứ hạng hạnh phúc.

\subsection{Quy trình tiền xử lý tổng quát}

\vspace{\baselineskip}
\begin{sodo}[H]
\centering
\fbox{\parbox[c][8.2cm][c]{0.88\textwidth}{\centering
\textbf{1. Đọc năm tệp CSV}\\[0.25cm]
$\downarrow$\\[0.2cm]
\textbf{2. Chuẩn hóa tên cột và tên quốc gia}\\[0.25cm]
$\downarrow$\\[0.2cm]
\textbf{3. Bổ sung Region và cột Year}\\[0.25cm]
$\downarrow$\\[0.2cm]
\textbf{4. Xử lý giá trị thiếu và loại trùng lặp}\\[0.25cm]
$\downarrow$\\[0.2cm]
\textbf{5. Nạp vào dashboard}
}}
\SodoCaption{Quy trình tiền xử lý dữ liệu}
\label{fig:quy-trinh-tien-xu-ly}
\end{sodo}
\vspace{\baselineskip}

\section{Kiến trúc và công nghệ triển khai}

\subsection{Kiến trúc tổng thể}

Ứng dụng được tổ chức theo mô hình nhiều trang của Streamlit. Tệp \textit{app.py} là trang chủ; thư mục \textit{pages} chứa các trang phân tích; \textit{dashboard\_utils.py} cung cấp hàm tải dữ liệu và giao diện chung; \textit{preprocess.py} thực hiện tiền xử lý dữ liệu; thư mục \textit{data} lưu các tệp nguồn theo năm.

\begin{table}[H]
\centering
\begin{tabular}{|p{4cm}|p{8.5cm}|}
\hline
\textbf{Thành phần} & \textbf{Vai trò} \\
\hline
app.py & Hiển thị trang chủ, thông tin tổng quan và điều hướng. \\
\hline
dashboard\_utils.py & Nạp dữ liệu, chuẩn hóa trường chung và áp dụng giao diện. \\
\hline
preprocess.py & Chuẩn hóa tên cột, bổ sung khu vực và xử lý giá trị thiếu. \\
\hline
pages/ & Chứa các trang Overview, Trends, Regional Analysis, Key Factors và Advanced Insights. \\
\hline
data/ & Lưu dữ liệu năm 2015--2019. \\
\hline
\end{tabular}
\caption{Vai trò của các thành phần trong hệ thống}
\label{tab:thanh-phan-he-thong}
\end{table}

\subsection{Công cụ và thư viện}

\begin{itemize}
    \item \textbf{Python}: ngôn ngữ lập trình chính.
    \item \textbf{Pandas}: nạp, nối, làm sạch, gom nhóm và tính toán dữ liệu.
    \item \textbf{Streamlit}: xây dựng giao diện dashboard nhiều trang và bộ lọc tương tác.
    \item \textbf{Plotly Express}: tạo bản đồ, biểu đồ đường, cột, vùng, phân tán và ma trận biểu đồ.
    \item \textbf{NumPy}: hỗ trợ xử lý dữ liệu số trong môi trường Python.
\end{itemize}

\section{Phương pháp thống kê và phân tích}

\subsection{Thống kê mô tả}

Ứng dụng tính giá trị trung bình của từng chỉ số theo năm bằng phương thức \textit{groupby} và \textit{mean}. Điểm hạnh phúc trung bình được dùng để theo dõi xu hướng chung, trong khi điểm trung bình theo khu vực được dùng để so sánh các vùng địa lý. Hệ thống cũng thống kê số quốc gia, số khu vực và số bản ghi trong phạm vi bộ lọc.

\subsection{Hệ số tương quan Pearson}

Mối liên hệ tuyến tính giữa hai biến được đo bằng hệ số tương quan Pearson. Với hai biến $X$ và $Y$, hệ số Pearson được tính bởi công thức:

\begin{equation}
r = \frac{\sum_{i=1}^{n}(x_i-\bar{x})(y_i-\bar{y})}
{\sqrt{\sum_{i=1}^{n}(x_i-\bar{x})^2}
\sqrt{\sum_{i=1}^{n}(y_i-\bar{y})^2}}
\label{pt:pearson}
\end{equation}

Trong đó:

\begin{itemize}
    \item $x_i$, $y_i$ là giá trị quan sát thứ $i$ của hai biến.
    \item $\bar{x}$, $\bar{y}$ là giá trị trung bình tương ứng.
    \item $r$ nằm trong khoảng từ $-1$ đến $1$.
\end{itemize}

Giá trị $r$ dương cho thấy hai biến có xu hướng tăng cùng chiều, giá trị $r$ âm cho thấy xu hướng ngược chiều và giá trị gần $0$ biểu thị liên hệ tuyến tính yếu. Tuy nhiên, tương quan không chứng minh quan hệ nhân quả.

\subsection{Biểu đồ phân tán và pairplot}

Biểu đồ phân tán được sử dụng để quan sát mối quan hệ giữa điểm hạnh phúc và một chỉ số khác. Pairplot cho phép hiển thị đồng thời nhiều cặp chỉ số; phần không thuộc đường chéo là biểu đồ phân tán, dữ liệu có thể được phân biệt theo khu vực hoặc năm bằng màu sắc.

\subsection{Gom nhóm và lọc dữ liệu}

Dashboard sử dụng ba loại bộ lọc chính:

\begin{itemize}
    \item Lọc theo năm báo cáo.
    \item Lọc theo khu vực địa lý hoặc quốc gia.
    \item Lọc theo khoảng điểm hạnh phúc.
\end{itemize}

Kết quả sau khi lọc được dùng trực tiếp để tính thống kê, vẽ biểu đồ và so sánh. Điều này giúp kết quả phân tích phản ánh đúng phạm vi mà người dùng đang quan sát.

\section{Thiết kế các chức năng phân tích}

\subsection{Trang tổng quan theo năm}

Trang Overview cho phép người dùng chọn năm, hiển thị bản đồ trực giao và lọc quốc gia theo khu vực. Khi chọn một khu vực trên bản đồ, bảng dữ liệu chỉ hiển thị các quốc gia thuộc khu vực đó.

\subsection{Trang xu hướng}

Trang Trends tổng hợp giá trị trung bình theo năm cho điểm hạnh phúc, GDP bình quân đầu người và tuổi thọ khỏe mạnh. Trang cũng cung cấp biểu đồ vùng điểm hạnh phúc và biểu đồ phân tán giữa điểm hạnh phúc với GDP hoặc tuổi thọ. Hệ số Pearson được tính từ dữ liệu sau khi loại các giá trị thiếu.

\subsection{Trang phân tích khu vực}

Trang Regional Analysis hỗ trợ lọc theo năm và nhiều khu vực. Hệ thống tính điểm hạnh phúc trung bình và số quốc gia cho từng khu vực, sau đó hiển thị bằng biểu đồ cột. Trang còn cho phép chọn ít nhất hai quốc gia để so sánh điểm hạnh phúc theo thời gian.

\subsection{Trang phân tích yếu tố}

Trang Key Factors tính ma trận tương quan Pearson cho điểm hạnh phúc, GDP, hỗ trợ xã hội, tuổi thọ khỏe mạnh, tự do lựa chọn, hào phóng và nhận thức về tham nhũng. Ma trận tương quan được hiển thị bằng heatmap. Trang cũng tính hệ số tương quan của từng yếu tố với điểm hạnh phúc và tạo pairplot từ các chỉ số được chọn.

\subsection{Trang khám phá nâng cao}

Trang Advanced Insights cung cấp bộ lọc phía thanh bên theo năm, khoảng điểm hạnh phúc và quốc gia. Trang được chia thành ba tab:

\begin{itemize}
    \item Tab khám phá hiển thị phân tán theo yếu tố và danh sách top 10 quốc gia.
    \item Tab so sánh theo dõi thay đổi điểm hạnh phúc của tối đa 8 quốc gia.
    \item Tab yếu tố phân tích scatter plot và tổng hợp tương quan của các yếu tố.
\end{itemize}

\section{Từ điển dữ liệu và quy tắc chuẩn hóa}

\subsection{Từ điển dữ liệu}

\begin{table}[H]
\centering
\begin{tabular}{|p{4.2cm}|p{2.4cm}|p{5.8cm}|}
\hline
\textbf{Tên trường} & \textbf{Kiểu dữ liệu} & \textbf{Ý nghĩa} \\
\hline
Country & Chuỗi & Tên quốc gia hoặc vùng lãnh thổ \\
\hline
Region & Chuỗi & Khu vực địa lý của quốc gia \\
\hline
Year & Số nguyên & Năm của báo cáo \\
\hline
Happiness Rank & Số nguyên & Thứ hạng hạnh phúc trong năm \\
\hline
Happiness Score & Số thực & Điểm hạnh phúc; dùng làm biến mục tiêu \\
\hline
GDP per capita & Số thực & Chỉ số GDP bình quân đầu người \\
\hline
Social support & Số thực & Mức độ hỗ trợ xã hội \\
\hline
Life Expectancy & Số thực & Tuổi thọ khỏe mạnh dự kiến \\
\hline
Freedom & Số thực & Tự do lựa chọn cuộc sống \\
\hline
Generosity & Số thực & Mức độ hào phóng \\
\hline
Perceptions of corruption & Số thực & Nhận thức về tham nhũng \\
\hline
\end{tabular}
\caption{Từ điển dữ liệu sau chuẩn hóa}
\label{tab:tu-dien-du-lieu}
\end{table}

Các tên trường trong Bảng \ref{tab:tu-dien-du-lieu} là tên nội bộ của ứng dụng. Chúng có thể khác với tên cột gốc trong từng tệp CSV. Việc định nghĩa một cấu trúc chung giúp các trang phân tích không phụ thuộc vào tên cột của từng năm.

\subsection{Quy tắc kiểm tra bản ghi}

Mỗi bản ghi cần có tên quốc gia, năm báo cáo và điểm hạnh phúc hợp lệ. Với bản ghi trùng quốc gia trong cùng một năm, chương trình giữ bản ghi có thứ hạng hạnh phúc hợp lý hoặc bản ghi đầu tiên sau khi dữ liệu đã được sắp xếp. Các trường số được ép kiểu bằng \textit{pd.to\_numeric}; giá trị không chuyển đổi được trở thành giá trị thiếu.

Các thao tác kiểm tra gồm:

\begin{itemize}
    \item Kiểm tra thiếu cột bắt buộc sau khi ánh xạ tên trường.
    \item Loại trùng quốc gia theo cặp \textit{Country, Year}.
    \item Kiểm tra điểm hạnh phúc và GDP không nằm ngoài phạm vi số hợp lệ.
    \item Loại hoặc bổ sung giá trị thiếu theo từng phạm vi phân tích.
    \item Kiểm tra kết quả gom nhóm không chứa nhóm rỗng trước khi vẽ biểu đồ.
\end{itemize}

\section{Thiết kế giao diện và luồng điều hướng}

\subsection{Trang chủ}

Trang chủ sử dụng các thẻ thống kê để giới thiệu giai đoạn dữ liệu, số quốc gia, điểm hạnh phúc trung bình và số khu vực. Phần điều hướng hiển thị năm trang phân tích chính nhằm người dùng có thể mở trực tiếp chức năng quan tâm.

\subsection{Thanh điều hướng và bố cục trang}

Ứng dụng sử dụng thanh điều hướng nhiều trang của Streamlit. Các bộ lọc chung được đặt ở thanh bên để người dùng không phải tìm kiếm trong từng khu vực nội dung. Nội dung chính được chia thành các hàng hoặc cột, hỗ trợ hiển thị đồng thời bảng số liệu và biểu đồ.

\subsection{Nguyên tắc hiển thị}

\begin{itemize}
    \item Giữ tên trục và tên đơn vị nhất quán giữa các trang.
    \item Đặt tên năm trên biểu đồ để dễ đối chiếu.
    \item Sử dụng thông điệp hoặc tiêu đề ngắn để mô tả chức năng của từng khu vực.
    \item Hiển thị giá trị trung bình, số lượng quan sát hoặc khoảng dữ liệu khi kết quả bị lọc.
    \item Tránh kết luận nhân quả nếu kết quả chỉ đến từ tương quan thống kê.
\end{itemize}

\section{Hướng dẫn cài đặt và chạy ứng dụng}

\subsection{Yêu cầu môi trường}

Môi trường triển khai cần Python và trình quản lý gói \textit{pip}. Các thư viện chính được khai báo trong tệp \textit{requirements.txt}, gồm Pandas, Streamlit, Plotly và các thư viện hỗ trợ cần thiết.

\subsection{Các bước chạy dự án}

\begin{enumerate}
    \item Tải mã nguồn từ kho GitHub hoặc sao chép thư mục dự án.
    \item Mở cửa sổ dòng lệnh tại thư mục gốc \textit{World-Happiness-Report\_Team11}.
    \item Cài thư viện bằng lệnh: \texttt{pip install -r requirements.txt}.
    \item Khởi chạy ứng dụng bằng lệnh: \texttt{streamlit run app.py}.
    \item Mở địa chỉ Streamlit hiển thị trên dòng lệnh để sử dụng dashboard.
\end{enumerate}

\subsection{Cấu trúc thư mục}

\begin{table}[H]
\centering
\begin{tabular}{|p{5.2cm}|p{7.2cm}|}
\hline
\textbf{Thư mục/tệp} & \textbf{Chức năng} \\
\hline
app.py & Cấu hình trang chủ và điều hướng \\
\hline
preprocess.py & Đọc, chuẩn hóa và làm sạch dữ liệu \\
\hline
dashboard\_utils.py & Hàm tải dữ liệu và cấu hình giao diện chung \\
\hline
pages/ & Các trang phân tích của dashboard \\
\hline
data/ & Tệp dữ liệu CSV năm 2015--2019 \\
\hline
report/ & Mã nguồn LaTeX của báo cáo \\
\hline
requirements.txt & Danh sách thư viện Python \\
\hline
\end{tabular}
\caption{Cấu trúc thư mục dự án}
\label{tab:cau-truc-thu-muc}
\end{table}

\section{Giả mã quy trình tiền xử lý}

Quy trình tiền xử lý có thể mô tả bằng các bước sau:

\begin{enumerate}
    \item Khởi tạo danh sách dữ liệu rỗng theo năm.
    \item Với từng tệp CSV, đọc dữ liệu và xác định năm tương ứng.
    \item Thay thế tên cột gốc bằng tên trường thống nhất.
    \item Ánh xạ tên quốc gia và bổ sung khu vực địa lý.
    \item Chuyển các trường số sang kiểu số và đánh dấu giá trị không hợp lệ.
    \item Loại bản ghi trùng quốc gia trong từng năm.
    \item Sắp xếp dữ liệu theo năm và thứ hạng hạnh phúc.
    \item Nối các năm thành một bộ dữ liệu chung cho dashboard.
\end{enumerate}

Chuỗi bước trên giúp phần phân tích chỉ làm việc với một cấu trúc dữ liệu nhất quán. Trong trường hợp tên cột hoặc tên quốc gia thay đổi, logic chuẩn hóa có thể được cập nhật tại module tiền xử lý thay vì sửa riêng từng trang.

\section{Quy trình phân tích dữ liệu trên dashboard}

\subsection{Thực hiện thống kê mô tả}

Với một bộ dữ liệu đã lọc, hệ thống dùng các phép gom nhóm theo năm hoặc khu vực và tính trung bình của trường cần phân tích. Số quan sát được lấy bằng số dòng dữ liệu sau lọc. Đối với top quốc gia, dữ liệu được sắp xếp giảm dần theo điểm hạnh phúc rồi lấy số dòng theo giới hạn hiển thị.

\subsection{Thực hiện tương quan}

Trước khi tính tương quan, hệ thống loại các cặp giá trị thiếu trong hai cột được chọn. Sau đó, \textit{corr()} hoặc công thức Pearson tương ứng được áp dụng để tính hệ số. Kết quả được làm tròn trước khi hiển thị nhằm giúp người dùng đọc dễ hơn.

\subsection{Cập nhật biểu đồ}

Mỗi khi người dùng thay đổi bộ lọc, Streamlit chạy lại phần cần thiết của trang. Dữ liệu trước tiên được lọc theo năm, khu vực hoặc khoảng điểm. Tiếp theo, biểu đồ được xây dựng lại từ dữ liệu đã lọc. Cách này giúp thống kê mô tả và biểu đồ quan sát cùng một phạm vi dữ liệu.

\section{Biện luận về độ tin cậy của kết quả}

Độ tin cậy của kết quả phụ thuộc vào chất lượng dữ liệu gốc, tính đầy đủ của trường cần phân tích và quy mô quan sát. Khi số quan sát lớn, trung bình khu vực và tương quan có xu hướng ổn định hơn. Khi một khu vực hoặc một quốc gia chỉ có ít quan sát, kết luận không nên được mở rộng ra toàn dân số.

Ngoài ra, dữ liệu World Happiness Report phản ánh các thang đo tổng hợp tại thời điểm khảo sát. Các biến có thể liên quan với nhau và chịu tác động của nhiều điều kiện kinh tế -- xã hội. Vì vậy, kết quả tương quan chỉ mô tả mối liên hệ tuyến tính trong tập dữ liệu chứ không xác định nguyên nhân.

\section{Các trường hợp dữ liệu đặc biệt}

Trong quá trình tổng hợp, dữ liệu có thể gặp tên quốc gia thay đổi, thiếu giá trị hoặc trùng bản ghi. Những trường hợp này được xử lý theo nguyên tắc tối thiểu làm sai dữ liệu gốc. Tên bí danh được dùng để nối khu vực với tên quốc gia, giá trị không hợp lệ được chuyển thành giá trị thiếu và bản ghi trùng được loại theo khóa quốc gia -- năm.

Nếu thiếu hoàn toàn một cột ở một năm, trang phân tích vẫn có thể sử dụng các cột còn có. Khi thiếu dữ liệu ở biến đầu vào của tương quan, cặp thiếu sẽ bị loại thay vì thay số thiếu bằng điểm hạnh phúc hoặc trung bình của khu vực.
